\chapter*{Вступ}
\addcontentsline{toc}{chapter}{Вступ}

% ============================================================
%  ВСТУП -- 11 канонічних блоків української докторської.
%  Блоки з похідним змістом (мета, об'єкт, предмет, новизна,
%  завдання) сформульовано; решта -- структурований каркас
%  під прозу. Наукова новизна узгоджена з 8-розділовою структурою.
% ============================================================

\section*{Актуальність теми}

% Каркас аргументації (розкрити прозою):
%  -- масштаб приватної цифрової комунікації та наскрізного шифрування;
%  -- зростання ризикових інформаційних впливів (зокрема для дітей);
%  -- суперечність між серверним контролем відкритого тексту та приватністю;
%  -- слабкість класифікації одного повідомлення без контексту;
%  -- постквантовий перехід і довготривала конфіденційність;
%  -- потреба контрольованого перегляду без відкритого тексту на сервері.
Сучасні захищені системи обміну повідомленнями з наскрізним шифруванням створюють принципову суперечність між приватністю комунікації та потребою захищати особу, насамперед дитину або іншу вразливу особу, від ризикових інформаційних впливів. \emph{[Розкрити актуальність прозою на основі Розділу~1.]}

\section*{Зв'язок роботи з науковими програмами, планами, темами}

\emph{[Зазначити НДР, держбюджетні/госпдоговірні теми, номери державної реєстрації, у межах яких виконано роботу, та роль здобувача в них.]}

\section*{Мета і завдання дослідження}

\textbf{Метою роботи} є розроблення методології, методів та інформаційної технології захисту контекстної приватності особи у захищених цифрових комунікаційних середовищах, що забезпечують оцінювання інформаційного ризику та прийняття захисних рішень без розкриття відкритого тексту повідомлень серверу, спираючись на постквантово-посилену автентифікацію, гібридне наскрізне шифрування та серверно-сліпий канал нагляду.

\medskip
\noindent Для досягнення мети поставлено такі \textbf{завдання}:

\begin{enumerate}
  \item проаналізувати стан проблеми захисту контекстної приватності в захищених системах обміну повідомленнями, обмеження серверного контролю за наскрізного шифрування та недостатність позаконтекстної класифікації повідомлень, виявити наукову прогалину (розділ~1);
  \item розробити концепцію інформаційної технології на основі середовища виконання політик AURA, узгоджену модель загроз та явні межі тверджень для всіх рівнів захисту (розділ~2);
  \item розробити метод контекстної оцінки інформаційного ризику, що враховує контакт, історію розмови, часову динаміку, повторюваність та латентні стани маніпуляції (розділ~3);
  \item розробити метод перетворення структурованого сигналу ризику в захисну дію з урахуванням політик, опрацюванням невизначеності та приватним аудитом (розділ~4);
  \item удосконалити криптографічну основу середовища виконання політик: метод постквантово-посиленої OPAQUE-автентифікації (розділ~5) та метод гібридного постквантового наскрізного шифрування (розділ~6);
  \item розробити метод серверно-сліпого каналу нагляду для передавання захисних сигналів пристрою опікуна без розкриття відкритого тексту серверу (розділ~7);
  \item реалізувати інформаційну технологію AURA/Ecliptix та виконати її експериментальну перевірку й оцінювання з порогами допуску, з визначенням обмежень і меж застосування (розділ~8).
\end{enumerate}

\section*{Об'єкт і предмет дослідження}

\textbf{Об'єкт дослідження} -- процеси інформаційного впливу, автентифікації, захищеного обміну повідомленнями та контекстної безпеки в цифрових комунікаційних середовищах.

\medskip
\noindent\textbf{Предмет дослідження} -- моделі, методи та програмні механізми захисту контекстної приватності особи, що охоплюють оцінювання ризику інформаційного впливу, приватне опрацювання захисних рішень, постквантово-посилену автентифікацію, гібридне наскрізне шифрування і серверно-сліпий канал нагляду.

\section*{Методи дослідження}

% Перелік прив'язати до конкретних розділів.
У роботі використано: методи системного аналізу та моделювання загроз (розділи~1--2); методи теорії інформаційної безпеки та криптографічного аналізу, зокрема формальну верифікацію протоколів засобами Tamarin та ProVerif (розділи~5--6); методи контекстного аналізу даних і прийняття рішень за умов невизначеності (розділи~3--4); методи інженерії програмного забезпечення, тестування та відтворюваної експериментальної перевірки з порогами допуску (розділ~8).

\section*{Наукова новизна одержаних результатів}

% Узгоджено з 8-розділовою структурою. Формули новизни -- за ДСТУ-стилем.
\begin{enumerate}
  \item \textbf{вперше запропоновано} архітектуру захисту контекстної приватності особи на основі середовища виконання політик AURA, у якій оцінювання ризику, захисні рішення з урахуванням політик, постквантово-посилена автентифікація, гібридне наскрізне шифрування та серверно-сліпий канал нагляду розглядаються як взаємопов'язані рівні однієї інформаційної технології, що, \emph{на відміну від наявних підходів}, не потребує перенесення відкритого тексту повідомлень на сервер (розділ~2);

  \item \textbf{вперше розроблено} метод контекстної оцінки інформаційного ризику в захищеному комунікаційному середовищі, у якому повідомлення розглядається як подія в контексті контакту, історії розмови, часової динаміки, повторюваності впливу та латентних станів маніпуляції, \emph{що дозволяє} перейти від класифікації окремого повідомлення до структурованого сигналу ризику (розділ~3);

  \item \textbf{удосконалено} метод прийняття рішень контекстної безпеки шляхом переходу від мітки ризику до простору захисних дій з урахуванням політик (дозвіл, попередження, ескалація, карантин, блокування, контрольований перегляд) з опрацюванням невизначеності та профілю захисту особи (розділ~4);

  \item \textbf{удосконалено} метод постквантово-посиленої парольної автентифікації на основі OPAQUE шляхом гібридного DH/KEM-узгодження сесії та явного операційного розмежування секрету \texttt{oprf\_seed}, \emph{що дозволяє} зберегти захист від перебору словника поза взаємодією із сервером за компрометації бази даних у межах визначеної моделі (розділ~5);

  \item \textbf{удосконалено} метод гібридного постквантового наскрізного шифрування шляхом KEM-підсилення кроку «храповика», ротації ключів метаданих та чіткого розмежування меж формального доведення, тестування й реалізаційної перевірки (розділ~6);

  \item \textbf{вперше запропоновано} метод серверно-сліпого каналу нагляду, у якому згода, контрольована пристроєм дитини, активний зв'язок між дитиною та опікуном і виділення парного ключа для кожного пристрою дають змогу передавати захисні сигнали пристрою опікуна без розкриття їхнього відкритого тексту серверу (розділ~7);

  \item \textbf{удосконалено} модель меж компрометації середовища виконання політик захисту контекстної приватності та методологію його приватного аудиту й оцінювання з порогами допуску шляхом розмежування сценаріїв компрометації і поєднання кураторських та синтетичних корпусів зі статусами \texttt{PASS}, \texttt{FAIL}, \texttt{INSUFFICIENT\_SUPPORT}, \texttt{BLOCKED} (розділи~2, 8);

  \item \textbf{розроблено} інформаційну технологію AURA/Ecliptix, яка реалізує запропоновані методи у вигляді програмного комплексу на основі Rust, iOS-застосунку та серверної частини з відтворюваною перевіркою й приватним аудитом (розділ~8).
\end{enumerate}

\section*{Практичне значення одержаних результатів}

\emph{[Розкрити: придатність методів і технології для захищених систем обміну повідомленнями, зокрема сценаріїв захисту дітей; впровадження в AURA/Ecliptix; відтворювані матеріали перевірки, тести, звіти допуску до випуску; акти впровадження за наявності.]}

\section*{Особистий внесок здобувача}

\emph{[Усі основні результати отримано здобувачем особисто. Для робіт у співавторстві -- конкретизувати внесок здобувача по кожній публікації.]}

\section*{Апробація результатів дисертації}

\emph{[Перелічити конференції, семінари, доповіді, де оприлюднено результати.]}

\section*{Публікації}

\emph{[Кількість і перелік публікацій за темою -- див. розгорнутий список у анотації. Зазначити статті у фахових виданнях та виданнях, проіндексованих у Scopus/Web of Science, і праці апробаційного характеру.]}

\section*{Структура та обсяг дисертації}

Дисертація складається зі вступу, восьми розділів, висновків, списку використаних джерел та додатків. \emph{[Уточнити загальний обсяг, кількість рисунків, таблиць, найменувань у списку джерел після завершення тексту.]}
